\section{Expand--Accumulate over a prime field}
\label{sec:prime-field-ea}

This section extends the enumerator proof to the weighted prime-field model
from Section~\ref{sec:prime-field-expander}.  Fix a prime $p$, put $s:=p-1$,
and measure the Hamming weight of a vector by its number of nonzero field
symbols.  The expander support is sampled as before, but every sampled edge
has an independent uniform label in $\Fbb_p^*$.  The accumulator remains the
prefix-sum map over $\Fbb_p$.

The random labels serve the same purpose as the nonzero diagonal weighters in
the generalized BAA construction: they make the law of a fixed message depend
only on its support.  Here they must be attached to expander edges, rather than
only after the expander, because a post-expander scaling does not change
whether an addition-only coordinate sum vanishes.

\subsection{The two-state accumulator kernel}

Suppose an accumulator input symbol is zero with probability $1-q$ and is
uniform on $\Fbb_p^*$ with total probability $q$.  If $z$ marks a nonzero
output symbol, the zero/nonzero state classes have transition matrix
\begin{equation}
 T^{(p)}_q(z):=
 \begin{pmatrix}
  1-q&qz\\
  q/s&(1-q/s)z
 \end{pmatrix}.
 \label{eq:prime-ea-weight-matrix}
\end{equation}
From a nonzero state, exactly one of the $s$ possible nonzero increments
returns the accumulator to zero.  All other increments preserve a nonzero
state.  This proves that the two classes are an exact lumping, not an
approximation.

For $q\in[0,s/p]$ and $z\in(0,1]$, define
\begin{equation}
\begin{split}
 \lambda_p(q,z):=\frac12\bigg(&1-q+(1-q/s)z\\
 &+\sqrt{\bigl(1-q-(1-q/s)z\bigr)^2+4q^2z/s}\bigg)
\end{split}
\label{eq:prime-ea-spectral-radius}
\end{equation}
and
\begin{equation}
 J^{(p)}_\delta(q)
 :=\sup_{0<z\leq1}
   \{\delta\ln z-\ln\lambda_p(q,z)\}.
 \label{eq:prime-ea-tail-rate}
\end{equation}

\begin{lemma}[Prime-field accumulator tail]
\label{lem:prime-ea-spectral-tail}
Let $H$ be the output weight of the $\Fbb_p$ accumulator on $n$ independent
inputs with the law above.  For $0<q\leq s/p$ and
$0<\delta<s/p$,
\begin{equation}
 \Pr[H\leq\lfloor\delta n\rfloor]
 \leq\sqrt p\exp\bigl(-nJ^{(p)}_\delta(q)\bigr).
 \label{eq:prime-ea-spectral-tail}
\end{equation}
\end{lemma}

\begin{proof}
Let $D_z:=\operatorname{diag}(1,\sqrt{sz})$.  The similar matrix
\[
 D_zT^{(p)}_q(z)D_z^{-1}
 =\begin{pmatrix}
   1-q&q\sqrt{z/s}\\
   q\sqrt{z/s}&(1-q/s)z
  \end{pmatrix}
\]
is symmetric and positive semidefinite for $q\leq s/p$.  Its largest
eigenvalue is \eqref{eq:prime-ea-spectral-radius}.  With $\ev_0$ selecting
state zero,
\[
 \Ebb[z^H]
 =\ev_0^{\mathsf T}T^{(p)}_q(z)^n\boldsymbol 1
 \leq\sqrt{1+sz}\,\lambda_p(q,z)^n
 \leq\sqrt p\,\lambda_p(q,z)^n.
\]
Markov's inequality applied to $z^H$, followed by optimization over
$z\in(0,1]$, proves the claim.
\end{proof}

For the labeled Bernoulli expander, let $q^{(p)}_r$ be
\eqref{eq:prime-bernoulli-activation}.  There are
\begin{equation}
 N_{p,k,r}:=\binom{k}{r}s^{r-1}
 \label{eq:prime-projective-support-count}
\end{equation}
one-dimensional message subspaces with support size $r$.  Counting these
projective messages, instead of all $s$ nonzero scalar multiples of each one,
gives the exact first-moment unit for the minimum-distance event.
For one fixed support of size $r$, write
\begin{equation}
 \tau^{(p)}_r(L)
 :=\sum_{h=0}^L[z^h]\,
   \ev_0^{\mathsf T}T^{(p)}_{q^{(p)}_r}(z)^n\boldsymbol 1
 \label{eq:prime-ea-line-tail}
\end{equation}
for the bad-output probability of any one projective message line on that
support.

\begin{theorem}[Labeled Bernoulli EA first moment]
\label{thm:prime-ea-first-moment}
Let $\B$ be the labeled Bernoulli expander of density $\theta$, let $\A$ be
the accumulator over $\Fbb_p$, and set $\G:=\B\A$.  The expected number of
one-dimensional message subspaces whose nonzero codewords have weight at most
$L$ is
\begin{equation}
 \boxed{
 \mu^{(p)}_{\mathrm{EA}}(L)
 =\sum_{r=1}^k N_{p,k,r}
   \sum_{h=0}^L[z^h]\,
   \ev_0^{\mathsf T}
   T^{(p)}_{q^{(p)}_r}(z)^n\boldsymbol 1.
 }
 \label{eq:prime-ea-first-moment}
\end{equation}
Consequently,
\begin{equation}
 \Pr[d_{\min}(\G)\leq L]
 \leq \mu^{(p)}_{\mathrm{EA}}(L).
 \label{eq:prime-ea-first-moment-failure}
\end{equation}
If $\theta\leq(p-1)/p$, then also
\begin{equation}
 \mu^{(p)}_{\mathrm{EA}}(L)
 \leq \sqrt p\sum_{r=1}^k N_{p,k,r}
       \exp\bigl(-nJ^{(p)}_{L/n}(q^{(p)}_r)\bigr).
 \label{eq:prime-ea-spectral-first-moment}
\end{equation}
\end{theorem}

\begin{proof}
Scalar multiples have the same output support, so it suffices to count one
representative of each message line.  Lemma~\ref{lem:prime-bernoulli-expander}
makes the expanded coordinates independent with nonzero probability
$q^{(p)}_r$.  Matrix \eqref{eq:prime-ea-weight-matrix} then gives the exact
output-weight generating function.  Sum over support sizes and output weights.
Markov's inequality bounds the bad-code event by the expected number of bad
message lines, and Lemma~\ref{lem:prime-ea-spectral-tail} gives the final
inequality.
\end{proof}

The preceding expectation treats all projective lines separately.  A first
improvement is to group the whole projective family carried by one message
support before taking the union bound.  This is the simplest form of the
projective-family truncation used in recent large-field RAA
analyses~\cite{akhiani26,khabbazian26}.

\begin{corollary}[Support-grouped Bernoulli bound]
\label{cor:prime-ea-support-grouped}
Under the hypotheses of Theorem~\ref{thm:prime-ea-first-moment},
\begin{equation}
 \boxed{
 \Pr[d_{\min}(\G)\leq L]
 \leq
 \sum_{r=1}^k\binom{k}{r}
 \min\!\left\{1,s^{r-1}\tau^{(p)}_r(L)\right\}.
 }
 \label{eq:prime-ea-support-grouped}
\end{equation}
If $\theta\leq s/p$, then also
\begin{equation}
 \Pr[d_{\min}(\G)\leq L]
 \leq
 \sum_{r=1}^k\binom{k}{r}
 \min\!\left\{1,
  s^{r-1}\sqrt p\exp\bigl(-nJ^{(p)}_{L/n}(q^{(p)}_r)\bigr)
 \right\}.
 \label{eq:prime-ea-support-grouped-spectral}
\end{equation}
\end{corollary}

\begin{proof}
Fix a support $S\subseteq[k]$ of size $r$.  It carries $s^{r-1}$
projective message lines.  Independent uniform edge labels make the bad-output
probability of every such line equal to $\tau^{(p)}_r(L)$.  The probability
that at least one line on $S$ is bad is therefore at most both $1$ and
$s^{r-1}\tau^{(p)}_r(L)$.  Sum this bound over the $\binom{k}{r}$ choices of
$S$.  Lemma~\ref{lem:prime-ea-spectral-tail} gives the second claim.
\end{proof}

Support grouping alone does not use the accumulator equations shared by the
projective family.  The following conditional bound does.  Fix an unlabeled
expander incidence pattern $\Gamma\subseteq[k]\times[n]$.  For
$U\subseteq[n]$, put $W=[n]\setminus U=\{w_1<\cdots<w_m\}$, set $w_0=0$, and
define the zero-terminated gaps
\begin{equation}
 I_a:=\{w_{a-1}+1,\ldots,w_a\},\qquad 1\leq a\leq m.
 \label{eq:prime-ea-zero-gaps}
\end{equation}
The trailing coordinates after $w_m$ impose no zero equation.  For a message
support $S$, let
\begin{equation}
 c_\Gamma(S,U)
 :=\left|\left\{a:\Gamma\cap(S\times I_a)\ne\emptyset\right\}\right|
 \label{eq:prime-ea-gap-rank}
\end{equation}
be the number of zero-terminated gaps hit by an active edge.

\begin{lemma}[Projective surplus cancellation for labeled EA]
\label{lem:prime-ea-projective-surplus}
Fix $\Gamma$, a message support $S$ of size $r\geq1$, and an output container
$U$.  Conditional on $\Gamma$, sample every edge label independently and
uniformly from $\Fbb_p^*$.  Then
\begin{equation}
 \Pr_{\mathrm{labels}}\!\left[
  \exists\xv:\operatorname{supp}(\xv)=S,\
  \operatorname{supp}(\xv\B_\Gamma\A)\subseteq U
 \right]
 \leq
 \boxed{\min\left\{1,s^{r-1-c_\Gamma(S,U)}\right\}.}
 \label{eq:prime-ea-projective-surplus}
\end{equation}
\end{lemma}

\begin{proof}
Fix one projective message $\xv$ with support $S$.  Requiring the accumulator
output to vanish at $w_a$ gives the equation
\[
 \sum_{(i,j)\in\Gamma:\ i\in S,\ j\leq w_a}x_i\lambda_{i,j}=0.
\]
Regard these $m$ equations as a homogeneous system in the active edge labels.
An edge ending in $I_a$ contributes, up to the nonzero column scale $x_i$,
the suffix vector that starts at equation $a$.  The distinct nonzero suffix
vectors are linearly independent.  The system therefore has rank exactly
$c_\Gamma(S,U)$, including all expander collisions.

If the active incidence contains $M$ edges, a rank-$c$ linear system has at
most $s^{M-c}$ solutions in $(\Fbb_p^*)^M$: choose its nonpivot variables,
after which the pivot variables are uniquely determined when they remain
nonzero.  The probability for the fixed projective message is at most
$s^{-c}$.  There are $s^{r-1}$ projective messages with exact support $S$.
Union-bound over them and cap the result at one.
\end{proof}

Consequently, for either the Bernoulli or regular unlabeled incidence law,
\begin{equation}
 \Pr[d_{\min}(\G)\leq L]
 \leq
 \sum_{r=1}^k\ \sum_{\substack{S\subseteq[k]\\|S|=r}}
 \ \sum_{\substack{U\subseteq[n]\\|U|=L}}
 \Ebb_\Gamma\!\left[
  \min\left\{1,s^{r-1-c_\Gamma(S,U)}\right\}
 \right].
 \label{eq:prime-ea-projective-master}
\end{equation}
Every output of weight at most $L$ is contained in at least one $L$-set $U$,
so this is a union over supports and output containers, not over field-valued
words or accumulator traces.

For Bernoulli incidence, the expectation in
\eqref{eq:prime-ea-projective-master} has an exact bivariate enumerator.  Put
$m=n-L$, $a_r=(1-\theta)^r$, and
\begin{equation}
 F_r(z,v)
 :=\frac{v}{1-z}+\frac{(1-v)a_r}{1-a_rz}.
 \label{eq:prime-ea-gap-enumerator}
\end{equation}
Then
\begin{equation}
 \boxed{
 \Pr[d_{\min}(\G)\leq L]
 \leq
 \sum_{r=1}^k\binom{k}{r}
 \sum_{c=0}^{m}
 [z^Lv^c]\frac{F_r(z,v)^m}{1-z}
 \min\left\{1,s^{r-1-c}\right\}.
 }
 \label{eq:prime-ea-projective-gap-enumerator}
\end{equation}
Indeed, write the $L$ coordinates of $U$ as $u_a\geq0$ coordinates before
the $a$th zero, plus a trailing part.  A gap of length $u_a+1$ is missed by
all edges leaving $S$ with probability $a_r^{u_a+1}$.  Summing its missed and
hit cases over $u_a$ gives $F_r$; the final factor $(1-z)^{-1}$ counts the
trailing part.  Thus the coefficient in
\eqref{eq:prime-ea-projective-gap-enumerator} sums the incidence probability
of every output container with exactly $c$ hit gaps.

Equation~\eqref{eq:prime-ea-projective-gap-enumerator} is rigorous but drops
the requirement that the coordinates inside $U$ are nonzero.  This loss is
severe for a sparse expander.  A message with $M$ active edges can change the
accumulator state at no more than $M$ coordinates, whereas the relaxed
container sum counts arbitrary subsets of $U$.  The RAA repair therefore has
two indispensable parts: the surplus-cancellation factor above and a placement
enumerator for the exact zero/nonzero accumulator trace.  The next EA bound
must insert $s^{-(c-(r-1))_+}$ into the existing sparse run enumerator rather
than use the container coefficient alone.

That insertion has a two-state form for Bernoulli incidence.  Let
\begin{equation}
 u_r:=1-(1-\theta)^r
 \label{eq:prime-ea-occupied-probability}
\end{equation}
be the probability that at least one edge leaves a fixed support of size $r$
and ends at a specified expanded coordinate.  Define
\begin{equation}
 \mathbf K_{u}(z,v):=
 \begin{pmatrix}
  1-u+uv&uz\\
  uv&z
 \end{pmatrix}.
 \label{eq:prime-ea-projective-trace-matrix}
\end{equation}
The variable $z$ marks nonzero accumulator outputs.  The variable $v$ marks
occupied coordinates at which the accumulator state is zero.

\begin{theorem}[Bernoulli trace-placement bound]
\label{thm:prime-ea-projective-trace}
For the labeled Bernoulli EA ensemble,
\begin{equation}
 \boxed{
 \Pr[d_{\min}(\G)\leq L]
 \leq
 \sum_{r=1}^k\binom{k}{r}
 \sum_{h=0}^L\sum_{a=0}^n
 [z^hv^a]\,
 \ev_0^{\mathsf T}\mathbf K_{u_r}(z,v)^n\boldsymbol 1
 \min\left\{1,s^{r-1-a}\right\}.
 }
 \label{eq:prime-ea-projective-trace}
\end{equation}
\end{theorem}

\begin{proof}
Fix a message support $S$ of size $r$.  Call an expanded coordinate
\emph{occupied} when at least one edge from $S$ ends there.  The occupancy
indicators are independent, and each has probability $u_r$.

Fix the occupied coordinates and a zero/nonzero accumulator trace.  At an
unoccupied coordinate, the accumulator state remains unchanged.  At an
occupied coordinate, the trace specifies whether the new state is zero or
nonzero.  Matrix \eqref{eq:prime-ea-projective-trace-matrix} enumerates these
choices and their output weights.

Suppose the trace has $a$ occupied coordinates whose new state is zero.  The
corresponding prefix equations in the active edge labels have rank $a$.
Indeed, subtracting consecutive equations isolates a nonempty set of fresh
edge variables.  Lemma~\ref{lem:prime-ea-projective-surplus} therefore bounds
the union over projective messages on $S$ by
$\min\{1,s^{r-1-a}\}$.  Sum over traces and over the $\binom{k}{r}$ supports.
\end{proof}

The exact coefficient bound has a convenient two-piece Chernoff relaxation.
For $z,v\in(0,1]$, define
\begin{align}
 A_r(z,v)
 &:=z^{-L}v^{-(r-1)}
   \ev_0^{\mathsf T}\mathbf K_{u_r}(z,v)^n\boldsymbol 1,
 \label{eq:prime-ea-projective-trace-structural}\\
 B_r(z,v)
 &:=s^{-1}z^{-L}v^{-r}
   \ev_0^{\mathsf T}\mathbf K_{u_r}(z,v)^n\boldsymbol 1.
 \label{eq:prime-ea-projective-trace-field}
\end{align}
Then
\begin{equation}
 \Pr[d_{\min}(\G)\leq L]
 \leq
 \sum_{r=1}^k\binom{k}{r}
 \left(
  \inf_{0<z,v\leq1}A_r(z,v)
  +\inf_{\substack{0<z\leq1\\s^{-1}\leq v\leq1}}B_r(z,v)
 \right).
 \label{eq:prime-ea-projective-trace-chernoff}
\end{equation}
The first term covers traces with $a\leq r-1$.  For $a\geq r$ and
$v\geq s^{-1}$,
\[
 s^{r-1-a}\leq s^{-1}v^{a-r},
\]
which gives the second term.  This split prevents the field-weighted branch
from reintroducing the projective factor on traces without surplus equations.

At rate $1/2$, the trace bound gives promising large-field parameters.  Take
$p=2^{127}-1$, $n=2{,}097{,}100$, and $k=1{,}048{,}550$.  The $p$-ary GV root
is $0.492127392425$.  The following rows report floating-point support grids.
\begin{center}
\begin{tabular}{c|rrrl}
 GV fraction&$L$&expected row weight&largest sampled term&support\\\hline
 $0.995$&$1{,}026{,}880$&$143$&$-78.1186$ bits&$r=1$\\
 floored GV&$1{,}032{,}040$&$194$&$-73.2977$ bits&$r=k$
\end{tabular}
\end{center}
These rows are diagnostics, not interval certificates.  They do not control
the unsampled support intervals.  For small fields, the ordinary per-line
enumerator can be stronger; the final finite proof should take the smaller
bound for each support size.

\subsection{Asymptotic distance}

Define the $p$-ary Hamming-shell entropy in nats by
\begin{equation}
 \mathsf h_p(x):=\mathsf h(x)+x\ln(p-1)
 \label{eq:prime-shell-entropy}
\end{equation}
and the sparse accumulator rate by
\begin{equation}
 I_{p,\delta}
 :=\left(\sqrt{1-\delta}
          -\sqrt{\frac{\delta}{p-1}}\right)^2.
 \label{eq:prime-ea-sparse-rate}
\end{equation}

\begin{lemma}[Prime-field endpoint exponents]
\label{lem:prime-ea-endpoints}
For $0<\delta<s/p$,
\begin{align}
 \liminf_{q\downarrow0}\frac{J^{(p)}_\delta(q)}q
 &\geq I_{p,\delta},
 \label{eq:prime-ea-sparse-limit}\\
 J^{(p)}_\delta(s/p)
 &=\ln p-\mathsf h_p(\delta).
 \label{eq:prime-ea-dense-rate}
\end{align}
\end{lemma}

\begin{proof}
For the sparse endpoint, set $z=e^{-\vartheta q}$.  Expanding
\eqref{eq:prime-ea-spectral-radius} at $q=0$ reduces the coefficient of $q$
in \eqref{eq:prime-ea-tail-rate} to
\[
 \frac{1+\vartheta+1/s}{2}
 -\frac12\sqrt{(\vartheta+1/s-1)^2+4/s}
 -\delta\vartheta.
\]
Its supremum over $\vartheta\geq0$ is
$(\sqrt{1-\delta}-\sqrt{\delta/s})^2$.

At $q=s/p$, the input is uniform on $\Fbb_p$ and
$\lambda_p(s/p,z)=(1+sz)/p$.  The optimizer is
$z=\delta/(s(1-\delta))$.  Substitution gives
\eqref{eq:prime-ea-dense-rate}.
\end{proof}

\begin{theorem}[Prime-field Expand--Accumulate distance]
\label{thm:prime-ea-linear-distance}
Fix a prime $p$.  Let $k/n\to R\in(0,1)$ and sample the labeled Bernoulli
expander with
\begin{equation}
 \theta_n=C\frac{\ln n}{n}.
 \label{eq:prime-ea-density}
\end{equation}
If $0<\delta<(p-1)/p$ and
\begin{align}
 C I_{p,\delta}&>1,
 \label{eq:prime-ea-sparse-condition}\\
 \mathsf h_p(\delta)&<(1-R)\ln p,
 \label{eq:prime-ea-gv-condition}
\end{align}
then
\begin{equation}
 \Pr\left[
  \min_{\xv\in\Fbb_p^k\setminus\{0\}}
  \hw{\xv\B\A}\leq\lfloor\delta n\rfloor
 \right]=o(1).
 \label{eq:prime-ea-linear-distance-conclusion}
\end{equation}
Thus the labeled prime-field ensemble reaches every distance strictly below
the $p$-ary Gilbert--Varshamov bound, subject only to the sparse density
condition \eqref{eq:prime-ea-sparse-condition}.
\end{theorem}

\begin{proof}
Use the three message-support ranges from the proof of
Theorem~\ref{thm:ea-linear-distance}.  For $r\theta_n=o(1)$,
\eqref{eq:prime-bernoulli-activation} gives
$q^{(p)}_r=r\theta_n(1-o(1))$.  Lemma~\ref{lem:prime-ea-endpoints} and
\eqref{eq:prime-ea-sparse-condition} therefore make the support-$r$ tail at
most $n^{-(1+\epsilon)r}$ for some $\epsilon>0$.  The projective support count
satisfies $N_{p,k,r}\leq(sk)^r$, so the sparse sum vanishes.

Once $r\theta_n$ is bounded below, $q^{(p)}_r$ is bounded away from zero.
Continuity and positivity of $J^{(p)}_\delta$ control supports up to
$r=\eta n$.  The number of such projective messages has exponent at most
$R\mathsf h(\eta/R)+\eta\ln s+o(1)$, which can be made arbitrarily small by
choosing $\eta>0$ small.

For $r>\eta n$, equation \eqref{eq:prime-bernoulli-activation} gives
$q^{(p)}_r=s/p-o(1)$ uniformly.  The total number of projective messages is at
most $p^k/s$.  By \eqref{eq:prime-ea-dense-rate}, their total contribution has
exponent
\[
 R\ln p+\mathsf h_p(\delta)-\ln p+o(1),
\]
which is negative by \eqref{eq:prime-ea-gv-condition}.  All three ranges
vanish in \eqref{eq:prime-ea-spectral-first-moment}.
\end{proof}

At rate $R=1/2$, the following values show the asymptotic trend.  The column
$\delta_{\mathrm{GV}}$ solves
$\mathsf h_p(\delta)=\tfrac12\ln p$, and $C_{\min}$ is the limiting sparse
threshold $1/I_{p,\delta_{\mathrm{GV}}}$.
\begin{center}
\begin{tabular}{c|cc}
 $p$&$\delta_{\mathrm{GV}}$&$C_{\min}$\\\hline
 $2$&$0.110028$&$2.67272$\\
 $3$&$0.159462$&$2.48437$\\
 $5$&$0.209867$&$2.29681$\\
 $7$&$0.237240$&$2.19795$\\
 $17$&$0.292842$&$2.00832$
\end{tabular}
\end{center}
The field increases the GV distance and, for these small primes, decreases the
required Bernoulli row-weight constant.

The quantifier ``fix a prime $p$'' is essential.  At a concrete length, the
projective count \eqref{eq:prime-projective-support-count} contributes
$(r-1)\ln(p-1)$ for support size $r$.  This term is constant in the
fixed-field asymptotic proof, but it is material when $p$ is a
cryptographic-size prime and $n$ is about $2^{20}$.  If
$\ln p/\ln n\to\alpha>0$, the sparse counting argument requires, at minimum,
a strengthened condition of the form $C I_{p,\delta}>1+\alpha$ in support
ranges where $r$ grows.  The theorem above therefore justifies small fixed
prime fields; it does not by itself justify a 128-bit field with the same
degree.

\subsection{Exact labeled regular composition}

The odd-prime regular ensemble has a qualitative advantage over the binary
one.  A binary even-weight message has even parity in every region, so the
accumulator returns to zero at every regional boundary.  Independent nonzero
edge labels remove that deterministic constraint.  For example, two active
rows have zero regional sum only when their two labels satisfy one scalar
equation, an event of probability $1/(p-1)$ rather than certainty.

The exact regional transfer remains two-state.  Let
\begin{equation}
 \mathbf U_p(X,Y):=
 \begin{pmatrix}
  1&sXY\\
  X&(1+(s-1)X)Y
 \end{pmatrix}.
 \label{eq:prime-ea-regular-bivariate-matrix}
\end{equation}
The exponent of $X$ marks nonzero inputs and the exponent of $Y$ marks
nonzero outputs.  For a uniform $p$-ary shell of weight $u$, define
\begin{align}
 \mathbf S^{(p),\mathrm{acc}}_{\ell,u}(Y)
 &:=\frac{[X^u]\mathbf U_p(X,Y)^\ell}
          {\binom{\ell}{u}s^u},
 \label{eq:prime-ea-regular-slice-matrix}\\
 \mathbf R^{(p),\mathrm{acc}}_{\ell,r}(Y)
 &:=\sum_{u=0}^\ell \rho^{(p,\ell)}_{r,u}
       \mathbf S^{(p),\mathrm{acc}}_{\ell,u}(Y).
 \label{eq:prime-ea-regular-region-matrix}
\end{align}

\begin{theorem}[Exact labeled regular EA first moment]
\label{thm:prime-ea-regular-first-moment}
Assume $n=d\ell$ and sample the labeled region-stratified left-regular
expander.  The expected number of bad projective message lines is
\begin{equation}
 \boxed{
 \mu^{(p),\mathrm{reg}}_{\mathrm{EA}}(L)
 =\sum_{r=1}^k N_{p,k,r}\sum_{h=0}^L[Y^h]\,
   \ev_0^{\mathsf T}
   \mathbf R^{(p),\mathrm{acc}}_{\ell,r}(Y)^d\boldsymbol 1.
 }
 \label{eq:prime-ea-regular-first-moment}
\end{equation}
It upper-bounds the probability of minimum distance at most $L$.
Moreover, the Poissonized bound is
\begin{equation}
 \Ebb_{\B}[z^{\hw{\xv\B\A}}]
 \leq\gamma_r^d\,
 \ev_0^{\mathsf T}
 T^{(p)}_{q^{(p),\mathrm{reg}}_r}(z)^n\boldsymbol 1,
 \label{eq:prime-ea-regular-poisson-mgf}
\end{equation}
and the dense bound is
\begin{equation}
 \Pr[\hw{\xv\B\A}\leq L]
 \leq V_{p,n}(L)p^{-n}
       (A_{p,\ell,r}+B_{p,\ell,r})^d,
 \quad
 V_{p,n}(L):=\sum_{h=0}^L\binom nhs^h.
 \label{eq:prime-ea-regular-dense-bound}
\end{equation}
\end{theorem}

\begin{proof}
Lemma~\ref{lem:prime-regular-shell-law} gives an independent mixture of
uniform $p$-ary shells in every region.  Matrix
\eqref{eq:prime-ea-regular-bivariate-matrix} counts all input values and keeps
the accumulator's boundary state, so coefficient extraction and normalization
give \eqref{eq:prime-ea-regular-first-moment}.  Equation
\eqref{eq:prime-regular-poisson-comparison} applied to the nonnegative output
moment gives \eqref{eq:prime-ea-regular-poisson-mgf}.  Finally, the accumulator
is invertible.  Its output ball has $V_{p,n}(L)$ preimages, and
Lemma~\ref{lem:prime-regular-point-mass} bounds the probability of each one.
\end{proof}

For the regular ensemble, define the per-line tail
\begin{equation}
 \tau^{(p),\mathrm{reg}}_r(L)
 :=\sum_{h=0}^L[Y^h]\,
   \ev_0^{\mathsf T}
   \mathbf R^{(p),\mathrm{acc}}_{\ell,r}(Y)^d\boldsymbol 1.
 \label{eq:prime-ea-regular-line-tail}
\end{equation}
Exactly the same support grouping gives
\begin{equation}
 \boxed{
 \Pr[d_{\min}(\G)\leq L]
 \leq
 \sum_{r=1}^k\binom{k}{r}
 \min\!\left\{1,
   s^{r-1}\tau^{(p),\mathrm{reg}}_r(L)
 \right\}.
 }
 \label{eq:prime-ea-regular-support-grouped}
\end{equation}
The same formula remains valid after replacing
$\tau^{(p),\mathrm{reg}}_r(L)$ by the optimized Poissonized or dense upper
bound from Theorem~\ref{thm:prime-ea-regular-first-moment}.  Thus the cap is
applied after union-bounding the projective lines on one support, but before
union-bounding the supports themselves.

The cap in \eqref{eq:prime-ea-regular-support-grouped} can remove a large
artificial contribution when the per-support line union already exceeds one.
Lemma~\ref{lem:prime-ea-projective-surplus} is sharper: it charges a factor
$s^{-1}$ only for each hit-gap equation beyond the first $r-1$.  It applies to
regular incidence without change.  We next perform the required incidence
average while retaining the exact accumulator trace.

For $t\geq0$, let $\omega^{(\ell)}_{t,u}$ be the probability that $t$
independent uniform draws from $[\ell]$ occupy exactly $u$ coordinates.  This
law is specified by $\omega^{(\ell)}_{0,0}=1$ and
\begin{equation}
 \omega^{(\ell)}_{t+1,u}
 =\frac{u}{\ell}\omega^{(\ell)}_{t,u}
  +\frac{\ell-u+1}{\ell}\omega^{(\ell)}_{t,u-1},
 \label{eq:prime-ea-regular-occupancy-law}
\end{equation}
where terms outside their natural range are zero.  Define the empty- and
occupied-coordinate trace matrices
\begin{equation}
 \mathbf Q_0(z):=
 \begin{pmatrix}1&0\\0&z\end{pmatrix},
 \qquad
 \mathbf Q_1(z,v):=
 \begin{pmatrix}v&z\\v&z\end{pmatrix}.
 \label{eq:prime-ea-regular-trace-inputs}
\end{equation}
Here $z$ marks a nonzero accumulator state and $v$ marks an occupied
coordinate whose new state is zero.  Put
\begin{align}
 \mathbf S^{\mathrm{tr}}_{\ell,u}(z,v)
 &:=\binom{\ell}{u}^{-1}[X^u]
   \bigl(\mathbf Q_0(z)+X\mathbf Q_1(z,v)\bigr)^\ell,
 \label{eq:prime-ea-regular-trace-slice}\\
 \mathbf R^{\mathrm{tr}}_{\ell,r}(z,v)
 &:=\sum_{u=0}^{\min\{r,\ell\}}
   \omega^{(\ell)}_{r,u}\mathbf S^{\mathrm{tr}}_{\ell,u}(z,v).
 \label{eq:prime-ea-regular-trace-region}
\end{align}

\begin{theorem}[Regular trace-placement bound]
\label{thm:prime-ea-regular-projective-trace}
For the labeled region-stratified left-regular EA ensemble with $n=d\ell$,
\begin{equation}
 \boxed{
 \Pr[d_{\min}(\G)\leq L]
 \leq
 \sum_{r=1}^k\binom{k}{r}
 \sum_{h=0}^L\sum_{a=0}^n[z^hv^a]\,
 \ev_0^{\mathsf T}
 \mathbf R^{\mathrm{tr}}_{\ell,r}(z,v)^d\boldsymbol 1
 \min\{1,s^{r-1-a}\}.
 }
 \label{eq:prime-ea-regular-projective-trace}
\end{equation}
\end{theorem}

\begin{proof}
Fix a message support $S$ of size $r$.  In one region, its $r$ rows make
independent uniform endpoint choices.  Their number of occupied coordinates
has law \eqref{eq:prime-ea-regular-occupancy-law}; conditional on that number,
the occupied set is uniform.  Equations
\eqref{eq:prime-ea-regular-trace-inputs}--
\eqref{eq:prime-ea-regular-trace-region} therefore enumerate the exact
incidence probability and every compatible zero/nonzero accumulator trace.
The $d$ regional endpoint choices are independent, while the two-state matrix
product carries the accumulator state across regional boundaries.

If a trace has $a$ occupied coordinates whose new state is zero, its $a$
prefix equations have rank $a$: subtracting consecutive equations isolates a
nonempty set of fresh edge-label variables.  Lemma~
\ref{lem:prime-ea-projective-surplus} bounds the union over the projective
messages on $S$ by $\min\{1,s^{r-1-a}\}$.  Summing the traces and supports
gives \eqref{eq:prime-ea-regular-projective-trace}.
\end{proof}

The exact coefficient sum admits the same two-piece Chernoff relaxation as
the Bernoulli bound.  For $z,v\in(0,1]$, put
\begin{align}
 A^{\mathrm{reg}}_r(z,v)
 &:=z^{-L}v^{-(r-1)}\ev_0^{\mathsf T}
   \mathbf R^{\mathrm{tr}}_{\ell,r}(z,v)^d\boldsymbol 1,
 \label{eq:prime-ea-regular-trace-structural}\\
 B^{\mathrm{reg}}_r(z,v)
 &:=s^{-1}z^{-L}v^{-r}\ev_0^{\mathsf T}
   \mathbf R^{\mathrm{tr}}_{\ell,r}(z,v)^d\boldsymbol 1.
 \label{eq:prime-ea-regular-trace-field}
\end{align}
Then the $r$th summand in
\eqref{eq:prime-ea-regular-projective-trace} is at most
\begin{equation}
 \binom{k}{r}\left(
  \inf_{0<z,v\leq1}A^{\mathrm{reg}}_r(z,v)
  +\inf_{\substack{0<z\leq1\\s^{-1}\leq v\leq1}}
   B^{\mathrm{reg}}_r(z,v)
 \right).
 \label{eq:prime-ea-regular-trace-chernoff}
\end{equation}

Two comparison bounds make this expression practical away from small $r$.
Let $\bar u_r=1-e^{-r/\ell}$ and let $\gamma_r$ be the Poisson conditioning
factor in Lemma~\ref{lem:regular-poisson-comparison}.  In
\eqref{eq:prime-ea-regular-trace-structural} and
\eqref{eq:prime-ea-regular-trace-field}, one may replace
\begin{equation}
 \ev_0^{\mathsf T}\mathbf R^{\mathrm{tr}}_{\ell,r}(z,v)^d
 \boldsymbol 1
 \quad\hbox{by}\quad
 \gamma_r^d\ev_0^{\mathsf T}\mathbf K_{\bar u_r}(z,v)^n\boldsymbol 1.
 \label{eq:prime-ea-regular-trace-poisson}
\end{equation}
This follows by Poissonizing the $r$ endpoint draws independently in every
region and then conditioning each regional total to equal $r$.

The fixed Poisson mean is not necessary.  The exact regional transfer has the
exponential-generating-function representation
\begin{equation}
 \mathbf R^{\mathrm{tr}}_{\ell,r}(z,v)
 =\frac{r!}{\ell^r}[t^r]
  \bigl(\mathbf Q_0(z)+(e^t-1)\mathbf Q_1(z,v)\bigr)^\ell.
 \label{eq:prime-ea-regular-trace-egf}
\end{equation}
Hence, for every $t>0$, it is bounded entrywise by
\begin{equation}
 \frac{r!}{\ell^rt^r}
  \bigl(\mathbf Q_0(z)+(e^t-1)\mathbf Q_1(z,v)\bigr)^\ell.
 \label{eq:prime-ea-regular-trace-saddle}
\end{equation}
Substituting \eqref{eq:prime-ea-regular-trace-saddle} into
\eqref{eq:prime-ea-regular-trace-structural} and
\eqref{eq:prime-ea-regular-trace-field} gives a third Chernoff marker, which
can be optimized jointly with $z$ and $v$.  The choice $t=r/\ell$ recovers
\eqref{eq:prime-ea-regular-trace-poisson}.  Optimizing $t$ can therefore only
improve the fixed-mean Poisson comparison.

There is also a comparison without a conditioning penalty.  Set
$\beta_r=(1-1/\ell)^r$ and
\begin{equation}
 \mathbf D_r(z,v):=\beta_r\mathbf Q_0(z)+\mathbf Q_1(z,v).
 \label{eq:prime-ea-regular-trace-dense-matrix}
\end{equation}
For a prescribed set $E$ of empty coordinates, with $e_g$ of them in region
$g$, the regular incidence law satisfies
\[
 \Pr[E\text{ is empty}]
 =\prod_{g=1}^d(1-e_g/\ell)^r
 \leq \prod_{g=1}^d(1-1/\ell)^{re_g}
 =\beta_r^{|E|}.
\]
Consequently, the same replacement is valid with
\begin{equation}
 \ev_0^{\mathsf T}\mathbf D_r(z,v)^n\boldsymbol 1.
 \label{eq:prime-ea-regular-trace-dense}
\end{equation}
The dense matrix is not stochastic; it upper-bounds each exact empty-set
probability before summing the compatible traces.  The exact transfer is best
for the first few support sizes, the Poisson comparison covers intermediate
sizes, and \eqref{eq:prime-ea-regular-trace-dense} avoids the accumulated
conditioning cost for dense supports.

Floating-point scans suggest that the removal of the binary regional-reset
constraint is quantitatively useful.  At rate $1/2$ and length near $2^{21}$,
the following candidates retain $99.5\%$ of the corresponding $p$-ary GV
distance.
\begin{center}
\begin{tabular}{c|rrrr}
 $p$&$d$&$\ell$&$L$&$L/n$\\\hline
 $3$&$44$&$47{,}662$&$332{,}739$&$0.1586641$\\
 $5$&$34$&$61{,}680$&$437{,}916$&$0.2088178$
\end{tabular}
\end{center}
For the ternary row, exact regional terms through support $24$, sampled
Poisson terms above it, and sampled dense Fourier terms are all below
$2^{-23.4}$.  For the quinary row, exact terms through support $48$ are below
$2^{-24.1}$ and the sampled remaining bounds are below $2^{-41.4}$.  These
are parameter diagnostics, not complete interval certificates: the unsampled
support intervals and outward rounding remain to be checked.  Relative to the
certified binary regular degree $64$, the candidate degrees reduce the edge
count by $31.25\%$ and $46.875\%$, respectively.

The projective trace bound also makes a direct large-field comparison
possible.  For $p=2^{127}-1$ and rate $1/2$, a floating-point scan that uses
the exact regional transfer through $r=5$, the Poisson comparison at
intermediate supports, and the dense empty-set bound near $r=k$ gives
\begin{center}
\begin{tabular}{c|rrrrrl}
 target&$d$&$\ell$&$n$&$L$&largest sampled term\\\hline
 $0.995\,\delta_{\mathrm{GV}}$&$143$&$14{,}666$&$2{,}097{,}238$
   &$1{,}026{,}947$&$-300.7780$ bits at $r=6$\\
 $\lfloor n\delta_{\mathrm{GV}}\rfloor-1$&$193$&$10{,}866$
   &$2{,}097{,}138$&$1{,}032{,}058$&$-98.2633$ bits at $r=k$
\end{tabular}
\end{center}
These remain support-grid diagnostics, not interval certificates.  Regularity
makes the first few support terms extremely small; for example, in the first
row the exact $r=1$ term is below $2^{-4860}$.  It does not reduce the degree
at these two large-field targets, however, because the dense full-support
term determines the frontier and is nearly the same as for Bernoulli
incidence.

\subsection{Two-sided regular incidence}

The preceding dense obstruction comes from rare output coordinates that no
active row hits.  Left regularity does not prevent those coordinates.  A
balanced regional ensemble removes them at full message support while
retaining an exact coefficient transfer.

Assume rate $1/2$, let $d$ be even, and write
$n=d\ell$, $k=n/2$, and $q=k/\ell=d/2$.  In each region, sample a uniform
permutation of the $k$ input rows into $k$ labeled slots.  Partition those
slots into $\ell$ groups of size $q$, and connect each row to the output
coordinate of its group.  Sample the edge labels independently in
$\Fbb_p^*$.  Thus every input row has one edge in every region, and every
output coordinate has exactly $q$ incident edges.

For a fixed input support of size $r$, the active slots in each region form a
uniform $r$-subset of the $k$ slots.  Therefore its exact regional trace
matrix is
\begin{equation}
 \mathbf R^{\mathrm{bi}}_{\ell,q,r}(z,v)
 :=\binom{k}{r}^{-1}[X^r]
 \left(
  \mathbf Q_0(z)+\bigl((1+X)^q-1\bigr)\mathbf Q_1(z,v)
 \right)^\ell.
 \label{eq:prime-ea-biregular-trace-region}
\end{equation}

\begin{theorem}[Two-sided regular trace bound]
\label{thm:prime-ea-biregular-projective-trace}
For the balanced regional ensemble above,
\begin{equation}
 \Pr[d_{\min}(\G)\leq L]
 \leq
 \sum_{r=1}^k\binom{k}{r}
 \sum_{h=0}^L\sum_{a=0}^n[z^hv^a]\,
 \ev_0^{\mathsf T}
 \mathbf R^{\mathrm{bi}}_{\ell,q,r}(z,v)^d\boldsymbol 1
 \min\{1,s^{r-1-a}\}.
 \label{eq:prime-ea-biregular-projective-trace}
\end{equation}
For every $x>0$, the regional matrix is bounded entrywise by
\begin{equation}
 \mathbf R^{\mathrm{bi}}_{\ell,q,r}(z,v)
 \preceq
 \frac{x^{-r}}{\binom{k}{r}}
 \left(
  \mathbf Q_0(z)+\bigl((1+x)^q-1\bigr)\mathbf Q_1(z,v)
 \right)^\ell.
 \label{eq:prime-ea-biregular-trace-saddle}
\end{equation}
\end{theorem}

\begin{proof}
Fix the input support.  A regional permutation maps it to a uniform set of
$r$ slots.  An output coordinate is empty when none of its $q$ slots is
active.  If exactly $j\geq1$ slots are active, the trace matrix remains
$\mathbf Q_1$ and the slot choices contribute $\binom qjX^j$.  This gives
\eqref{eq:prime-ea-biregular-trace-region}.  The proof of Theorem~
\ref{thm:prime-ea-regular-projective-trace} then gives
\eqref{eq:prime-ea-biregular-projective-trace}.  Finally, all matrix
coefficients are nonnegative, so positive-marker coefficient extraction gives
\eqref{eq:prime-ea-biregular-trace-saddle}.
\end{proof}

The coefficient saddle also yields a finite block bound.  Fix $x,z,v>0$ on
an interval $I=[a,b]$ of message-support sizes.  After the union over input
supports, the only $r$-dependent prefactor is
\begin{equation}
 g_r:=\binom{k}{r}^{1-d}(x^dv)^{-r}.
 \label{eq:prime-ea-biregular-block-factor}
\end{equation}
The function $r\mapsto\log g_r$ is convex because
$r\mapsto\log\binom kr$ is concave and $1-d<0$.  Hence
\begin{equation}
 \sum_{r=a}^b g_r
 \leq |I|\max\{g_a,g_b\}
 \leq |I|(g_a+g_b).
 \label{eq:prime-ea-biregular-block-bound}
\end{equation}
The last expression avoids an interval comparison between $g_a$ and $g_b$.
It can be evaluated directly with outward-rounded arithmetic.  The structural
branch multiplies it by $v$; the field branch multiplies it by $s^{-1}$.

\begin{theorem}[Finite two-sided regular EA certificate]
\label{thm:prime-ea-biregular-finite-certificate}
Let $p=2^{127}-1$, and consider the two-sided regular ensemble with
\[
 d=30,\qquad q=15,\qquad \ell=69{,}905,\qquad
 n=2{,}097{,}150,\qquad k=1{,}048{,}575.
\]
For $L=1{,}032{,}064$,
\begin{equation}
 \boxed{
 \Pr[d_{\min}(\G)\leq L]
 <9.499\cdot10^{-10}<2^{-29.97}.
 }
 \label{eq:prime-ea-biregular-finite-certificate}
\end{equation}
The cutoff equals $\lfloor n\delta_{\mathrm{GV}}\rfloor$ for the rate-$1/2$
$p$-ary GV root $\delta_{\mathrm{GV}}=0.492127392425\ldots$.
\end{theorem}

The certificate evaluates every $1\leq r\leq64$ with the exact regional
matrix \eqref{eq:prime-ea-biregular-trace-region}.  Twelve applications of
\eqref{eq:prime-ea-biregular-block-bound} cover $65\leq r<k$, and the verifier
checks $r=k$ directly.  At 192-bit Arb precision, the three contributions are
bounded by
\[
 7.319\cdot10^{-30},\qquad
 9.499\cdot10^{-10},\qquad
 1.036\cdot10^{-75}.
\]
The verifier treats the floating-point optimizer's decimal markers as fixed
inputs.  Thus floating-point arithmetic affects efficiency, but not the
validity of \eqref{eq:prime-ea-biregular-finite-certificate}.  Degree $28$
has negative selected exact terms, but its present coefficient bound is
positive at the exact-to-saddle transition.

\subsection{Prime-field random convolution}

The accumulator certificate motivates the corresponding Expand--Convolute
question.  We randomize the feedback coefficients over the full field.  This
choice preserves the reduced state used in the binary proof.  Restricting
every coefficient to $\Fbb_p^*$ would instead make the transition law depend
on the support and values of the convolution state.

Fix a memory $m\geq1$.  Independently sample
$\boldsymbol\alpha_t\gets\Fbb_p^m$ for every $t\in[n]$.  Given
$\mathbf u\in\Fbb_p^n$, set $y_j:=0$ for $j\leq0$ and compute
\begin{equation}
 y_t:=u_t+\left\langle\boldsymbol\alpha_t,
       (y_{t-1},\ldots,y_{t-m})\right\rangle,
 \qquad t\in[n].
 \label{eq:prime-ec-recurrence}
\end{equation}
Write $\C_{\boldsymbol\alpha}$ for this nonwrapping convolution.  Every
realization is invertible because the coefficient of $u_t$ is one.

Index the transfer states by $0,\ldots,m$.  State $j<m$ records exactly $j$
trailing zero outputs, while state $m$ means that the convolution state is
zero.  Let $\mathbf e_m$ be the unit vector for state $m$, let
$\boldsymbol 1$ be the all-one vector, and put $s=p-1$.  The transfer below
counts constraints; its entries are not transition probabilities.  Define
the empty-coordinate matrix $\mathbf Q^{\mathrm{ec}}_0(z,v)$ and the
occupied-coordinate matrix $\mathbf Q^{\mathrm{ec}}_1(z,v)$ by the following
nonzero entries:
\begin{align}
 (\mathbf Q^{\mathrm{ec}}_b)_{j,0}
   &=z,
 &
 (\mathbf Q^{\mathrm{ec}}_b)_{j,j+1}
   &=v
 &&(b\in\{0,1\},\ 0\leq j<m),
 \label{eq:prime-ec-active-constraint-trace}\\
 (\mathbf Q^{\mathrm{ec}}_0)_{m,m}&=1,
 &
 (\mathbf Q^{\mathrm{ec}}_1)_{m,0}&=z,
 &
 (\mathbf Q^{\mathrm{ec}}_1)_{m,m}&=v.
 \label{eq:prime-ec-zero-constraint-trace}
\end{align}
The marker $z$ records a nonzero output.  The marker $v$ records an occupied
zero event that imposes an equation in fresh randomness.  Such an event
occurs either when a nonzero convolution state produces zero or when an
occupied expander coordinate produces zero in state $m$.

Use the two-sided regular expander from
Theorem~\ref{thm:prime-ea-biregular-projective-trace}.  For a message support
of size $r$, define its regional convolution trace by
\begin{equation}
 \mathbf R^{\mathrm{bi,ec}}_{\ell,q,r}(z,v)
 :=\binom{k}{r}^{-1}[X^r]
 \left(
 \mathbf Q^{\mathrm{ec}}_0(z,v)
  +\bigl((1+X)^q-1\bigr)\mathbf Q^{\mathrm{ec}}_1(z,v)
 \right)^\ell.
 \label{eq:prime-ec-biregular-region}
\end{equation}

\begin{theorem}[Two-sided regular prime-field EC constraint-trace bound]
\label{thm:prime-ec-biregular-projective-trace}
Independently sample the balanced labeled expander $\B$ and the convolution
$\C_{\boldsymbol\alpha}$ from \eqref{eq:prime-ec-recurrence}.  Set
$\G:=\B\C_{\boldsymbol\alpha}$.  Then
\begin{equation}
 \boxed{
 \Pr[d_{\min}(\G)\leq L]
 \leq
 \sum_{r=1}^k\binom{k}{r}
 \sum_{h=0}^L\sum_{a=0}^n[z^hv^a]\,
 \mathbf e_m^{\mathsf T}
 \mathbf R^{\mathrm{bi,ec}}_{\ell,q,r}(z,v)^d\boldsymbol 1
 \min\{1,s^{r-1-a}\}.
 }
 \label{eq:prime-ec-biregular-projective-trace}
\end{equation}
For every $x>0$, one may replace the regional matrix entrywise by
\begin{equation}
 \frac{x^{-r}}{\binom{k}{r}}
 \left(
  \mathbf Q^{\mathrm{ec}}_0(z)
  +\bigl((1+x)^q-1\bigr)\mathbf Q^{\mathrm{ec}}_1(z,v)
 \right)^\ell.
 \label{eq:prime-ec-biregular-saddle}
\end{equation}
\end{theorem}

\begin{proof}
Fix a message support of size $r$, one projective message line on that
support, and one zero--nonzero output trace.  Expose the randomness in time
order.  If the convolution state is nonzero and the next output is zero,
then \eqref{eq:prime-ec-recurrence} imposes one nontrivial affine equation on
the fresh vector $\boldsymbol\alpha_t$.  Its conditional probability is
$1/p\leq1/s$.  This is the transition marked by $v$ in
\eqref{eq:prime-ec-active-constraint-trace}.  We bound the complementary
nonzero event by one and mark its output weight by $z$.

If the convolution state is zero, the next output is the expander symbol.
An empty coordinate remains in state $m$.  At an occupied coordinate, a zero
symbol imposes one nontrivial linear equation on edge labels used only at
that coordinate.  Conditioning on all but one nonzero label bounds its
probability by $1/s$.  This is the transition marked by $v$ in
\eqref{eq:prime-ec-zero-constraint-trace}.  A nonzero symbol moves to state
zero and contributes $z$.

Every marked event uses a fresh feedback vector or fresh coordinate labels.
Consequently, a trace with $a$ marks has probability at most $s^{-a}$ for
the fixed projective line.  There are $s^{r-1}$ projective message lines on
the support.  The union bound, capped by one, is therefore
$\min\{1,s^{r-1-a}\}$.  Averaging the balanced incidence pattern gives
\eqref{eq:prime-ec-biregular-region}.  Matrix multiplication retains the
convolution state across regional boundaries.  Summing over traces, supports,
and output weights proves
\eqref{eq:prime-ec-biregular-projective-trace}.  Positive-marker coefficient
extraction proves \eqref{eq:prime-ec-biregular-saddle}.
\end{proof}

The trace above treats every occupied right coordinate alike.  The balanced
expander supplies one additional distinction.  A coordinate incident to
exactly one supported message symbol cannot vanish because both the message
symbol and its edge label are nonzero.  Define
\(\mathbf Q^{\mathrm{ec}}_{\mathrm{col}}:=
\mathbf Q^{\mathrm{ec}}_1\).  Define the singleton matrix
\(\mathbf Q^{\mathrm{ec}}_{\mathrm{sg}}\) to agree with
\(\mathbf Q^{\mathrm{ec}}_1\), except that
\begin{equation}
 (\mathbf Q^{\mathrm{ec}}_{\mathrm{sg}})_{m,m}:=0.
 \label{eq:prime-ec-singleton-matrix}
\end{equation}
Thus a singleton coordinate in state $m$ must produce a nonzero output.
For a message support of size $r$, put
\begin{equation}
 \widehat{\mathbf R}^{\mathrm{bi,ec}}_{\ell,q,r}(z,v)
 :=\binom{k}{r}^{-1}[X^r]
 \left(
  \mathbf Q^{\mathrm{ec}}_0
  +qX\mathbf Q^{\mathrm{ec}}_{\mathrm{sg}}
  +\bigl((1+X)^q-1-qX\bigr)
     \mathbf Q^{\mathrm{ec}}_{\mathrm{col}}
 \right)^\ell.
 \label{eq:prime-ec-singleton-region}
\end{equation}

\begin{corollary}[Singleton-refined EC trace]
\label{cor:prime-ec-singleton-trace}
Theorem~\ref{thm:prime-ec-biregular-projective-trace} remains valid when
\(\mathbf R^{\mathrm{bi,ec}}_{\ell,q,r}\) is replaced by
\(\widehat{\mathbf R}^{\mathrm{bi,ec}}_{\ell,q,r}\).
For positive coefficient extraction, define
\begin{equation}
 \widehat{\mathbf M}(x,z,v)
 :=\mathbf Q^{\mathrm{ec}}_0
   +qx\mathbf Q^{\mathrm{ec}}_{\mathrm{sg}}
   +\bigl((1+x)^q-1-qx\bigr)
      \mathbf Q^{\mathrm{ec}}_{\mathrm{col}}.
 \label{eq:prime-ec-singleton-saddle-matrix}
\end{equation}
\par\noindent
Then the entrywise saddle bound replaces the regional matrix by
\begin{equation*}
 \frac{x^{-r}}{\binom{k}{r}}\widehat{\mathbf M}(x,z,v)^\ell.
\end{equation*}
\end{corollary}

\begin{proof}
Fix one region and condition on its assignment of the $r$ supported rows to
the $k=q\ell$ slots.  A right coordinate receives zero, one, or at least two
supported slots.  These cases contribute the three terms in
\eqref{eq:prime-ec-singleton-region}.  In the singleton case, the expander
symbol is the product of two nonzero field elements.  Hence the zero
transition from state $m$ is impossible.  All remaining transitions use the
bounds from Theorem~\ref{thm:prime-ec-biregular-projective-trace}.
\end{proof}

The structural branch can still be loose at small support because its
equation marker also counts traces with $a\geq r$.  We isolate the traces
with $a<r$ before applying a marker.  For a fixed support of size $r$, call a
region \emph{singleton-free} if no right coordinate receives exactly one
supported slot.  Let $U_r$ be the number of occupied right coordinates after
a uniform $r$-subset of the $q\ell$ regional slots is selected, and define
\begin{equation}
 \beta_r:=\Pr[U_r\leq\lfloor r/2\rfloor].
 \label{eq:prime-ec-singleton-free-probability}
\end{equation}
A singleton-free region has at most $\lfloor r/2\rfloor$ occupied
coordinates.  Hence $\beta_r$ upper-bounds its probability.

\begin{lemma}[Singleton-free-region bound]
\label{lem:prime-ec-singleton-free-regions}
Fix a support size $r$ and consider traces with $a\leq r-1$ fresh equations.
Put
\begin{align}
 J_r&:=1+\left\lfloor\frac{r-1}{m}\right\rfloor,\\
 b_r&:=\left\lceil\frac{n-L-(r-1)}{\ell}\right\rceil-2J_r.
 \label{eq:prime-ec-required-singleton-free-regions}
\end{align}
If $b_r>0$, the contribution of these traces to
\eqref{eq:prime-ec-biregular-projective-trace} is at most
\begin{equation}
 \binom{k}{r}\binom{d}{b_r}\beta_r^{b_r}.
 \label{eq:prime-ec-singleton-free-region-bound}
\end{equation}
\end{lemma}

\begin{proof}
The convolution begins in state $m$.  Each later entry into state $m$
requires $m$ consecutive zero outputs from a nonzero state.  Those outputs
impose $m$ fresh equations.  A trace with $a\leq r-1$ therefore has at most
$J_r$ maximal intervals in state $m$.

A state-$m$ interval can meet at most two regions that contain a singleton
coordinate.  Otherwise the interval encounters a singleton coordinate in a
fully traversed region and ends there.  If $B$ regions are singleton-free,
the state-$m$ intervals contain at most $(2J_r+B)\ell$ coordinates.  The
remaining zero outputs occur outside state $m$ and impose fresh equations.
Thus a trace with output weight at most $L$ satisfies
\[
 n-L\leq (2J_r+B)\ell+a
       \leq (2J_r+B)\ell+r-1.
\]
It follows that $B\geq b_r$.  Regional permutations are independent, and
each region is singleton-free with probability at most $\beta_r$.  A union
bound over $b_r$ regions gives
\eqref{eq:prime-ec-singleton-free-region-bound}.
\end{proof}

For completeness, we record how the finite verifier avoids a support-by-support
union-bound loss.  For an integer interval $I\subseteq[1,k]$, define
\begin{align}
 \mathcal U_I^{\mathrm{str}}(z,v)
 &:={z^{-L}}\sum_{r\in I}\binom{k}{r}
 \mathbf e_m^{\mathsf T}
 \widehat{\mathbf R}^{\mathrm{bi,ec}}_{\ell,q,r}(z,v)^d\boldsymbol 1
 v^{-(r-1)},
 \label{eq:prime-ec-exact-band-structural}\\
 \mathcal U_I^{\mathrm{fld}}(z,v)
 &:={s^{-1}z^{-L}}\sum_{r\in I}\binom{k}{r}
 \mathbf e_m^{\mathsf T}
 \widehat{\mathbf R}^{\mathrm{bi,ec}}_{\ell,q,r}(z,v)^d\boldsymbol 1
 v^{-r}.
 \label{eq:prime-ec-exact-band-field}
\end{align}
The first expression is valid for $0<z,v\leq1$.  The second is valid for
$0<z\leq1$ and $s^{-1}\leq v\leq1$.  Thus, for independently selected
valid marker pairs $(z_{\mathrm{str}},v_{\mathrm{str}})$ and
$(z_{\mathrm{fld}},v_{\mathrm{fld}})$, the exact band bound is
\begin{equation}
 \mathcal U_I^{\mathrm{str}}(z_{\mathrm{str}},v_{\mathrm{str}})
 +\mathcal U_I^{\mathrm{fld}}(z_{\mathrm{fld}},v_{\mathrm{fld}}).
 \label{eq:prime-ec-exact-band}
\end{equation}
Indeed, the first summand dominates the traces with $a\leq r-1$, and the
second dominates those with $a\geq r$.  Each marker pair bounds its branch
for an entire finite support band.  The verifier constructs all normalized
regional coefficients in that band from one truncated matrix polynomial.

Larger intervals use convex endpoint domination.  Put
\begin{align}
 g_r(x,v)&:=\binom{k}{r}^{1-d}(x^dv)^{-r}.
\end{align}
For $I=[r_0,r_1]\subseteq[1,k-1]$, the structural and field contributions
are at most, respectively,
\begin{align}
 |I|vz^{-L}\mathbf e_m^{\mathsf T}\widehat{\mathbf M}(x,z,v)^n\boldsymbol 1
   \bigl(g_{r_0}(x,v)+g_{r_1}(x,v)\bigr),
 \label{eq:prime-ec-structural-block}\\
 \frac{|I|}{s}z^{-L}\mathbf e_m^{\mathsf T}\widehat{\mathbf M}(x,z,v)^n\boldsymbol 1
   \bigl(g_{r_0}(x,v)+g_{r_1}(x,v)\bigr).
 \label{eq:prime-ec-field-block}
\end{align}
The two lines may use different positive markers.  To justify the endpoint
step, observe that
$\ln g_r=(1-d)\ln\binom{k}{r}-r\ln(x^dv)$ is convex in $r$.

\begin{corollary}[Finite rate-half prime-field EC certificate]
\label{cor:prime-ec-biregular-finite-certificate}
Let $p=2^{127}-1$, $d=26$, $q=13$, $\ell=80{,}659$,
$n=2{,}097{,}134$, $k=1{,}048{,}567$, and $m=4$.  Let
$\delta_{\mathrm{GV}}$ solve
$\mathsf h_p(\delta_{\mathrm{GV}})=\tfrac12\ln p$, and set
$L=\lfloor n\delta_{\mathrm{GV}}\rfloor=1{,}032{,}057$.  Then
\begin{equation}
 \Pr[d_{\min}(\B\C_{\boldsymbol\alpha})\leq L]
 \leq 3.052485\cdot10^{-22}
 <2^{-71.4724}<2^{-20}.
 \label{eq:prime-ec-biregular-finite-certificate}
\end{equation}
\end{corollary}

\begin{proof}
The outward-rounded verifier applies \eqref{eq:prime-ec-exact-band} to eleven
intervals covering $[1,800]$.  Their sum is at most
$3.052485\cdot10^{-22}$.  It applies
\eqref{eq:prime-ec-structural-block} and
\eqref{eq:prime-ec-field-block} to $26$ intervals covering $[801,k-1]$.
Their sum is at most $1.005\cdot10^{-62}$.  Direct evaluation of the
full-support transfer adds less than $2.770\cdot10^{-42}$.  Summing the three
outward-rounded quantities proves
\eqref{eq:prime-ec-biregular-finite-certificate}.
\end{proof}

The certificate stores only fixed decimal markers and integer interval
endpoints.  Floating-point optimization is outside the trusted base.  The
Arb verifier recomputes every matrix, binomial coefficient, exponentiation,
and comparison at 256-bit precision.  From the \texttt{enumerator\_paper}
directory, it is run with
\begin{verbatim}
python scripts/prime_field_biregular_ec_singleton_certificate.py \
  verify CERTIFICATE
\end{verbatim}
Here \texttt{CERTIFICATE} denotes
\path{results/prime_field_biregular_ec_p127_rate_half_d26_m4_singleton_gv.json}.

The preceding certificate minimizes the simple work proxy $d+m$ but does not
minimize the expander degree.  Increasing the convolution memory by two permits
two fewer expander edges per input coordinate.

\begin{corollary}[Lower-degree rate-half prime-field EC certificate]
\label{cor:prime-ec-biregular-d24-m6-certificate}
Let $p=2^{127}-1$, $d=24$, $q=12$, $\ell=87{,}381$,
$n=2{,}097{,}144$, $k=1{,}048{,}572$, and $m=6$.  Let
$\delta_{\mathrm{GV}}$ solve
$\mathsf h_p(\delta_{\mathrm{GV}})=\tfrac12\ln p$, and set
$L=\lfloor n\delta_{\mathrm{GV}}\rfloor=1{,}032{,}062$.  Then
\begin{equation}
 \Pr[d_{\min}(\B\C_{\boldsymbol\alpha})\leq L]
 \leq 1.285949\cdot10^{-12}
 <2^{-39.5003}<2^{-20}.
 \label{eq:prime-ec-biregular-d24-m6-certificate}
\end{equation}
\end{corollary}

\begin{proof}
The verifier uses $17$ applications of \eqref{eq:prime-ec-exact-band} to cover
$[1,700]$.  Their sum is at most
$1.285949\cdot10^{-12}$.  The interval $[225,256]$ contributes the largest
term, which is at most $1.211048\cdot10^{-12}$.  The verifier applies
\eqref{eq:prime-ec-structural-block} and \eqref{eq:prime-ec-field-block} to
$30$ intervals covering $[701,k-1]$.  Their sum is at most
$1.392\cdot10^{-60}$.  Direct evaluation of the full-support transfer adds
less than $2.840\cdot10^{-39}$.  Summing these three outward-rounded
quantities proves \eqref{eq:prime-ec-biregular-d24-m6-certificate}.
\end{proof}

The lower-degree certificate is verified from the \texttt{enumerator\_paper}
directory with
\begin{verbatim}
python scripts/prime_field_biregular_ec_d24_m6_certificate.py \
  verify CERTIFICATE24
\end{verbatim}
The symbol \texttt{CERTIFICATE24} denotes the following file.
\begin{flushleft}\small
\path{results/prime_field_biregular_ec_p127_rate_half_d24_m6_singleton_gv.json}
\end{flushleft}

The preceding Chernoff split counts high-equation traces in its structural
branch.  Lemma~\ref{lem:prime-ec-singleton-free-regions} removes this slack
at small support and permits two fewer expander edges.

\begin{corollary}[Degree-22 rate-half prime-field EC certificate]
\label{cor:prime-ec-biregular-d22-m12-certificate}
Let $p=2^{127}-1$, $d=22$, $q=11$, $\ell=95{,}325$,
$n=2{,}097{,}150$, $k=1{,}048{,}575$, and $m=12$.  Let
$\delta_{\mathrm{GV}}$ solve
$\mathsf h_p(\delta_{\mathrm{GV}})=\tfrac12\ln p$, and set
$L=\lfloor n\delta_{\mathrm{GV}}\rfloor=1{,}032{,}064$.  Then
\begin{equation}
 \Pr[d_{\min}(\B\C_{\boldsymbol\alpha})\leq L]
 \leq 6.071992\cdot10^{-9}
 <2^{-27.2951}<2^{-20}.
 \label{eq:prime-ec-biregular-d22-m12-certificate}
\end{equation}
\end{corollary}

\begin{proof}
For supports $1\leq r\leq48$, the verifier applies
Lemma~\ref{lem:prime-ec-singleton-free-regions} to the branch $a<r$.  It
applies \eqref{eq:prime-ec-exact-band-field} to the branch $a\geq r$.
The verifier then applies \eqref{eq:prime-ec-exact-band} to $21$ intervals
covering $[49,1400]$.  The interval $[481,520]$ contributes the largest term,
which is at most $4.759588\cdot10^{-9}$.

The verifier applies \eqref{eq:prime-ec-structural-block} and
\eqref{eq:prime-ec-field-block} to $26$ intervals covering $[1401,k-1]$.
Their sum is at most $6.868\cdot10^{-98}$.  The full-support term is less than
$1.036\cdot10^{-75}$.  Summing the outward-rounded quantities proves
\eqref{eq:prime-ec-biregular-d22-m12-certificate}.
\end{proof}

The degree-$22$ certificate is verified with
\begin{verbatim}
python scripts/prime_field_biregular_ec_d22_m12_certificate.py \
  verify CERTIFICATE22
\end{verbatim}
where \texttt{CERTIFICATE22} denotes
\begin{flushleft}\small
\path{results/prime_field_biregular_ec_p127_rate_half_d22_m12_region_gv.json}.
\end{flushleft}

Table~\ref{tab:prime-ec-memory-degree-frontier} records the observed tradeoff
between convolution memory and expander degree.  The scan uses even $d$
because rate one half requires the right degree $q=d/2$.
\begin{table}[H]
\centering
\caption{Observed memory--degree frontier at rate one half, length near
$2^{21}$, and the floored $p$-ary GV cutoff.  Certified rows bound the
complete support range with 256-bit Arb arithmetic.  Diagnostic rows report
the best candidate found by the exact-term scan.}
\label{tab:prime-ec-memory-degree-frontier}
\begin{tabular}{c|c|c|c|l}
memory $m$ & left degree $d$ & right degree $q$ & $d+m$ & status \\ \hline
1  & 38 & 19 & 39 & diagnostic \\
2  & 30 & 15 & 32 & diagnostic; narrow margin \\
3  & 28 & 14 & 31 & certified; $27.65$ bits \\
4  & 26 & 13 & 30 & certified; $71.47$ bits \\
6  & 24 & 12 & 30 & certified; $39.50$ bits \\
12 & 22 & 11 & 34 & certified; $27.30$ bits
\end{tabular}
\end{table}

The middle of the table gives the smallest value of the work proxy $d+m$.
Memory four and memory six both give $d+m=30$.  The returns then diminish:
memory twelve reaches degree $22$ only after separating the low-equation
traces from the structural Chernoff bound.
It does not reach degree $20$ by memory $21=\lceil\log_2 n\rceil$.  At
$d=20$, $m=21$, and support $512$, the exact structural upper bound is
$2^{131.90}$.  A positive upper bound does not imply that the ensemble
fails; it identifies a parameter point not certified by this analysis.
